\section{总体分析}

四个子问题呈"由静到动、由单到联"的递进结构。问题一定标：完成统计与预测并给出可行基础调度，为后续三问提供统一的负荷—能源结算基线；问题二在无储能条件下释放"任务柔性格"，通过任务迁移改变逐时负荷以追逐低边际成本时段；问题三冻结负荷、释放"储能的灵活性格"，通过充放电实现新能源时序搬运与外送套利；问题四同时释放两类灵活性，探究二者的协同与替代关系，并通过情景与参数扫描刻画策略的鲁棒性与适用边界。

贯穿四个问题的主线是同一条能源结算模型：给定任意逐时设施负荷，先给出该区域可消纳/可外送/须购电的电力平衡，再核算成本、碳排与新能源利用率。该结算作为"黑箱"，使任务调度与储能优化两个问题解耦为"决策变量$ \to $负荷$ \to $结算价值"的复合结构，从而使问题二到问题四都能在相同评价口径下公平比较，避免因口径不一致导致的结论失真。

需要特别说明的数据事实：六个区域的可用新能源逐时序列完全相同（E/F 标签并不意味光伏/风电互补），基线运行记录中存在区域 F 在 2400 时刻的净负荷缺口（约 $31.6$ MW）以及部分区域终端 SOC 不守恒等现象；本文在建模中以物理守恒（电力平衡、SOC 递归与终端约束）重新结算，而非直接沿用基线记录中的购售电数值，以确保口径统一与可验证性。
